\begin{exercise}
Let $n$ be a fixed positive integer and let $a$ and $b$ be any two integers. Recall that $a$ is congruent
to $b$ modulo $n$ if $n$ divides $a-b$. Calculate:
\begin{enumerate*}[label=\alph*), itemjoin={\qquad}]
\item $17 \bmod 4$
\item $4 \bmod 17$
\item $(-2) \bmod 5$
\end{enumerate*}
\end{exercise}

\begin{exercise}
Solve the equations:
\begin{enumerate}[label=\alph*)]
\item $2x \equiv 1 \pmod 7$
\item $2x \equiv 1 \pmod 6$
\item $5x \equiv 1 \pmod 9$
\item $x^2 \equiv 3 \pmod{11}$
\item $2x+3 \equiv 0 \pmod 5$
\item $x+k \equiv 0 \pmod n$
\end{enumerate}
\end{exercise}

\begin{exercise}
Prove that for all $n,k,p\in \Z$ we have the following:
\begin{enumerate}[label=\alph*)]
\item $(n \bmod k) \bmod k = n \bmod k$
\item $(n+p) \bmod k = \big((n \bmod k) + (p \bmod k)\big) \bmod k$
\item $(np) \bmod k = \big((n \bmod k)\,(p \bmod k)\big) \bmod k$
\end{enumerate}
\end{exercise}

\begin{exercise}
Recall that for $a,b\in\N$ the operations of addition and multiplication modulo $n$ are defined by
$a +_n b = (a+b) \bmod n$ and $a \cdot_n b = (a\cdot b) \bmod n$. Prove that:
\begin{enumerate}[label=\alph*)]
\item addition modulo $n$ is associative, i.e.\ $(a +_n b) +_n c = a +_n (b +_n c)$ for any $a,b,c\in\N$;
\item multiplication modulo $n$ is associative, i.e.\ $(a \cdot_n b) \cdot_n c = a \cdot_n (b \cdot_n c)$
for any $a,b,c\in\N$;
\item multiplication modulo $n$ is distributive with respect to addition modulo $n$, i.e.\
$(a +_n b) \cdot_n c = (a \cdot_n c) +_n (b \cdot_n c)$ for any $a,b,c\in\N$.
\end{enumerate}
\end{exercise}

\begin{exercise}
State which of the following are groups:
\begin{enumerate*}[label=\alph*), itemjoin={\qquad}]
\item $(\Z_n, +_n)$
\item $(\Z_n\setminus\{0\}, \cdot_n)$
\item $(\{1,2,3,4,6,12\}, \gcd)$
\end{enumerate*}
\end{exercise}

\begin{exercise}
Show that $(2^X, \bigtriangleup)$ is a group, where $\bigtriangleup$ denotes the symmetric difference.
Construct the group table of $(2^{\{a,b,c\}}, \bigtriangleup)$.
\end{exercise}

\begin{exercise}
Show that $(\Z_n\setminus\{0\}, \cdot_n)$ is a group if and only if $n$ is a prime number.
\end{exercise}

\begin{exercise}
Verify whether $(\R_+, \bullet)$ is a group, where $a\bullet b = a^2 b^2$.
\end{exercise}

\begin{exercise}
Let $p\in\N$ be a prime number. Show that $(\Z_p, +_p, \cdot_p)$ is a field.
\end{exercise}

\begin{exercise}
Determine which of the following structures are fields. Do this in two steps: first verify whether the
set is a group with respect to the first operation, and then whether the set without the identity element
of the first operation is a group with respect to the second operation. Here $2^X$ carries the operations
$\cup$, $\cap$ and the symmetric difference $\bigtriangleup$, while $\R^\R$ denotes the set of all
functions $\R\to\R$, with pointwise addition $+$ and composition $\circ$.
\begin{enumerate*}[label=\alph*), itemjoin={\qquad}]
\item $(2^X, \cup, \cap)$
\item $(\R^\R, +, \circ)$
\item $(2^X, \cap, \cup)$
\item $(\R^\R, \circ, +)$
\item $(2^X, \cap, \bigtriangleup)$
\item $(2^X, \bigtriangleup, \cap)$
\item $(2^X, \cup, \bigtriangleup)$
\item $(2^X, \bigtriangleup, \cup)$
\end{enumerate*}
\end{exercise}

\begin{exercise}
Verify whether $(\R, \oplus, \odot)$ is a field, where $a \oplus b = a+b+1$ and $a\odot b = ab+a+b$.
\end{exercise}
